% SGA 3, Expose VI_A, Loop 1 English reconstruction.
% Scope: paragraph 5.2 in full; paragraph 5.3 is excluded.
% Authority: Polo--Gille Exp6A-13oct24.pdf, component-local page 31
% (combined-reader page 334; printed page 323).
% OCR and the Reinhold English Markdown are comparison witnesses only.

\medskip
\phantomsection\label{par:sga3-VIA-5.2}%
\noindent\textbf{5.2.---}
当 \(A\)-群同态 \(u:F\to G\) 为单态射时,
\hyperref[par:sga3-VIA-5.1]{5.1} 的结果也可以由第~IV~篇的结果得到.
事实上, 由 \hyperref[thm:sga3-VIA-3.2]{3.2}, 典范态射
\[
  p:G\longrightarrow F\backslash G
\]
忠实平坦且为开态射; 因而由\SGARefLink{sga3:iv:proposition:6:3:1}{第~IV~篇 6.3.1},
它是 \((\mathrm{fpqc})\) 拓扑下的覆盖态射,
于是可以应用\SGARefLink{sga3:iv:corollary:5:2:2}{第~IV~篇推论~5.2.2} 和 5.2.4.

特别地, 除 \hyperref[thm:sga3-VIA-3.2]{3.2} 的假设外,
再设 \(u\) 是正规子群 \(F\) 到 \(G\) 的包含.
则 \(F\backslash G\) 上存在唯一的 \(A\)-群结构, 使得典范态射
\[
  p:G\longrightarrow F\backslash G
\]
为 \(A\)-群同态.%
\footnote{\label{note:sga3-VIA-51}编者注: 以下这句话为增补内容.}
这就证明了\hyperref[concl:sga3-VIA-3.3.2-v]{定理~3.3.2(v)}.
% Independent prose review: source and translation compared paragraph by paragraph through this component; statements, hypotheses, proofs, notes and diagram captions covered; bibliography reserved for its assigned reviewer; source math and reference keys unchanged.
